\section{Manual de instalación y uso}\label{anx:manual}
\textbf{Requisitos:} Python $\ge$ 3.10, en Windows, Linux o macOS.
\begin{lstlisting}
git clone <URL-del-repositorio> && cd transporte-cusco-scheduling
python -m venv .venv
source .venv/bin/activate              # Windows: .venv\Scripts\activate
pip install -r requirements.txt
pip install -e .

python -m bus_sched build              # genera data/instances/*.json (35 rutas + DEMO)
python -m bus_sched profiles           # muestra los perfiles T0, T1 y T2
python -m bus_sched demo --traffic T2  # instancia manual con LS y LPT
python -m bus_sched run --route RTU-05 --alg ls  --traffic T1 --gantt
python -m bus_sched run --route RTU-05 --alg lpt --traffic T1
python experiments/run_experiments.py  # experimento E1
python experiments/build_excel.py      # Excel de la entrega
python experiments/export_latex.py     # tablas y figuras de este informe
python -m pytest -q                    # pruebas
\end{lstlisting}

\begin{table}[H]
\centering\small
\caption{Parámetros de la línea de comandos.}
\begin{tabularx}{\textwidth}{L{4.2cm}L{5.2cm}Y}
\toprule
Parámetro & Valores & Efecto \\
\midrule
\texttt{--route} / \texttt{--instance} & código de ruta, DEMO o archivo JSON & Instancia a resolver \\
\texttt{--alg} & \texttt{ls} · \texttt{lpt} & Algoritmo de la Entrega 1 \\
\texttt{--traffic} & T0 · T1 · T2 · archivo.json & Perfil de tráfico \\
\texttt{--gantt}, \texttt{--out} & --- & Diagrama de Gantt y ruta del CSV de salida \\
\bottomrule
\end{tabularx}
\end{table}

\section{Contribuciones y bitácora}\label{anx:contrib}
R = responsable, A = apoyo, V = revisión. La distribución se contrasta con el historial del repositorio.
\begin{table}[H]
\centering\small
\caption{Matriz de contribuciones.}
\begin{tabular}{lcccc}
\toprule
Componente & Aguilar & Choque & Huacani & Moreano \\
\midrule
Formulación y modelo de tráfico & R & V & V & A \\
Modelo de datos y línea de tiempo & R & A & V & \\
List Scheduling y LPT & A & R & V & \\
Validador, métricas y pruebas & A & A & R & A \\
Datos, Excel e instancias & V & & & R \\
Experimento E1 y gráficos & A & & V & R \\
Ejemplo manual & A & R & R & V \\
Informe & R & A & A & A \\
\bottomrule
\end{tabular}
\end{table}
Cada integrante mantiene una bitácora individual en \texttt{docs/bitacora/}, con una fila por sesión de trabajo: fecha, conocimiento trabajado, actividad, fuente, lo aprendido, dónde se aplicó (archivo o commit), dificultad y solución, y horas.

\section{Resultados por ruta}\label{anx:rutas}
$\rho$ = saturación de la flota; NA = vueltas no atendidas; $\Lmax$ en horas. Datos completos en \texttt{results/E1\_35\_rutas.csv} y en la hoja \emph{Resultados por ruta} del Excel.

{\footnotesize
\begin{longtable}{lrrrrrrrrrr}
\caption{Resultados por ruta (resaltadas: rutas con $\rho \geq 0.85$).}\label{tab:rutas}\\
\toprule
Ruta & Buses & Vueltas & $\rho$ & NA T0 & \makecell{LS\\NA T1} & \makecell{LS\\$L_{\max}$ T1} & \makecell{LPT\\NA T1} & \makecell{LPT\\$L_{\max}$ T1} & \makecell{LS\\NA T2} & \makecell{LPT\\NA T2} \\
\midrule\endfirsthead
\toprule
Ruta & Buses & Vueltas & $\rho$ & NA T0 & \makecell{LS\\NA T1} & \makecell{LS\\$L_{\max}$ T1} & \makecell{LPT\\NA T1} & \makecell{LPT\\$L_{\max}$ T1} & \makecell{LS\\NA T2} & \makecell{LPT\\NA T2} \\
\midrule\endhead
\bottomrule\endfoot
RTI-01 & 35 & 103 & 0.52 & 0 & 0 & 7.77 & 3 & 9.73 & 4 & 13 \\
\rowcolor{dorado!20}RTI-02 & 27 & 127 & 0.87 & 15 & 19 & 10.74 & 36 & 10.70 & 29 & 43 \\
RTI-04 & 42 & 144 & 0.55 & 1 & 2 & 8.64 & 4 & 9.04 & 5 & 16 \\
RTI-05 & 44 & 150 & 0.48 & 0 & 1 & 7.55 & 0 & 7.87 & 2 & 3 \\
RTI-06 & 32 & 144 & 0.56 & 0 & 3 & 8.53 & 0 & 8.67 & 4 & 8 \\
RTI-07 & 39 & 144 & 0.64 & 4 & 8 & 9.46 & 15 & 9.90 & 11 & 28 \\
RTI-08 & 20 & 108 & 0.61 & 0 & 1 & 9.23 & 1 & 9.31 & 3 & 11 \\
RTI-09 & 43 & 135 & 0.54 & 0 & 2 & 9.21 & 8 & 9.54 & 6 & 17 \\
RTI-10 & 37 & 166 & 0.62 & 1 & 4 & 9.48 & 11 & 9.73 & 7 & 19 \\
RTU-01 & 32 & 144 & 0.79 & 12 & 15 & 11.93 & 28 & 10.12 & 22 & 39 \\
RTU-02 & 29 & 117 & 0.63 & 2 & 4 & 8.91 & 12 & 8.95 & 8 & 18 \\
RTU-04 & 28 & 144 & 0.75 & 5 & 7 & 10.53 & 21 & 10.35 & 13 & 29 \\
RTU-05 & 24 & 144 & 0.82 & 7 & 10 & 11.50 & 28 & 9.91 & 18 & 36 \\
RTU-07 & 28 & 127 & 0.46 & 0 & 0 & 7.03 & 0 & 7.11 & 1 & 0 \\
RTU-08 & 26 & 114 & 0.55 & 0 & 2 & 8.64 & 2 & 8.77 & 3 & 6 \\
RTU-09 & 21 & 120 & 0.65 & 0 & 1 & 9.70 & 5 & 9.60 & 4 & 12 \\
RTU-10 & 16 & 103 & 0.61 & 0 & 0 & 9.11 & 2 & 9.26 & 1 & 6 \\
RTU-11A & 21 & 103 & 0.42 & 0 & 0 & 6.22 & 0 & 6.13 & 0 & 0 \\
RTU-11B & 22 & 103 & 0.40 & 0 & 0 & 6.06 & 0 & 6.16 & 0 & 0 \\
RTU-12 & 39 & 150 & 0.52 & 2 & 3 & 7.62 & 0 & 7.77 & 4 & 5 \\
RTU-13 & 37 & 144 & 0.65 & 3 & 8 & 9.32 & 16 & 11.65 & 11 & 26 \\
RTU-14 & 37 & 108 & 0.42 & 0 & 0 & 6.35 & 0 & 6.21 & 0 & 0 \\
RTU-15 & 30 & 127 & 0.58 & 1 & 4 & 9.21 & 3 & 9.58 & 5 & 11 \\
RTU-16 & 40 & 150 & 0.53 & 2 & 2 & 7.86 & 0 & 7.97 & 5 & 10 \\
RTU-17 & 28 & 127 & 0.57 & 0 & 2 & 8.72 & 1 & 8.95 & 3 & 7 \\
RTU-18 & 35 & 166 & 0.58 & 0 & 4 & 8.54 & 0 & 10.12 & 5 & 8 \\
RTU-19 & 34 & 144 & 0.70 & 6 & 10 & 9.51 & 21 & 9.44 & 14 & 31 \\
RTU-20 & 31 & 154 & 0.55 & 0 & 2 & 7.87 & 2 & 7.86 & 4 & 6 \\
RTU-21 & 26 & 166 & 0.59 & 0 & 0 & 8.86 & 0 & 9.00 & 0 & 6 \\
\rowcolor{dorado!20}RTU-22 & 19 & 127 & 1.02 & 21 & 25 & 12.94 & 39 & 12.83 & 37 & 52 \\
RTU-23 & 37 & 166 & 0.60 & 1 & 5 & 9.23 & 4 & 9.39 & 5 & 19 \\
RTU-24 & 35 & 117 & 0.58 & 2 & 4 & 9.42 & 7 & 9.70 & 7 & 22 \\
RTU-25 & 27 & 127 & 0.56 & 0 & 2 & 8.29 & 1 & 8.42 & 3 & 5 \\
RTU-26 & 33 & 154 & 0.66 & 2 & 5 & 9.79 & 16 & 10.00 & 9 & 23 \\
RTU-27 & 26 & 166 & 0.65 & 0 & 0 & 9.73 & 7 & 9.84 & 2 & 12 \\
\midrule
\textbf{Total} & 1080 & 4733 & & 87 & 155 & & 293 & & 255 & 547 \\
\end{longtable}}


